% Options for packages loaded elsewhere
\PassOptionsToPackage{unicode}{hyperref}
\PassOptionsToPackage{hyphens}{url}
\documentclass[
]{article}
\usepackage{xcolor}
\usepackage[margin=1in]{geometry}
\usepackage{amsmath,amssymb}
\setcounter{secnumdepth}{-\maxdimen} % remove section numbering
\usepackage{iftex}
\ifPDFTeX
  \usepackage[T1]{fontenc}
  \usepackage[utf8]{inputenc}
  \usepackage{textcomp} % provide euro and other symbols
\else % if luatex or xetex
  \usepackage{unicode-math} % this also loads fontspec
  \defaultfontfeatures{Scale=MatchLowercase}
  \defaultfontfeatures[\rmfamily]{Ligatures=TeX,Scale=1}
\fi
\usepackage{lmodern}
\ifPDFTeX\else
  % xetex/luatex font selection
\fi
% Use upquote if available, for straight quotes in verbatim environments
\IfFileExists{upquote.sty}{\usepackage{upquote}}{}
\IfFileExists{microtype.sty}{% use microtype if available
  \usepackage[]{microtype}
  \UseMicrotypeSet[protrusion]{basicmath} % disable protrusion for tt fonts
}{}
\makeatletter
\@ifundefined{KOMAClassName}{% if non-KOMA class
  \IfFileExists{parskip.sty}{%
    \usepackage{parskip}
  }{% else
    \setlength{\parindent}{0pt}
    \setlength{\parskip}{6pt plus 2pt minus 1pt}}
}{% if KOMA class
  \KOMAoptions{parskip=half}}
\makeatother
\usepackage{color}
\usepackage{fancyvrb}
\newcommand{\VerbBar}{|}
\newcommand{\VERB}{\Verb[commandchars=\\\{\}]}
\DefineVerbatimEnvironment{Highlighting}{Verbatim}{commandchars=\\\{\}}
% Add ',fontsize=\small' for more characters per line
\usepackage{framed}
\definecolor{shadecolor}{RGB}{248,248,248}
\newenvironment{Shaded}{\begin{snugshade}}{\end{snugshade}}
\newcommand{\AlertTok}[1]{\textcolor[rgb]{0.94,0.16,0.16}{#1}}
\newcommand{\AnnotationTok}[1]{\textcolor[rgb]{0.56,0.35,0.01}{\textbf{\textit{#1}}}}
\newcommand{\AttributeTok}[1]{\textcolor[rgb]{0.13,0.29,0.53}{#1}}
\newcommand{\BaseNTok}[1]{\textcolor[rgb]{0.00,0.00,0.81}{#1}}
\newcommand{\BuiltInTok}[1]{#1}
\newcommand{\CharTok}[1]{\textcolor[rgb]{0.31,0.60,0.02}{#1}}
\newcommand{\CommentTok}[1]{\textcolor[rgb]{0.56,0.35,0.01}{\textit{#1}}}
\newcommand{\CommentVarTok}[1]{\textcolor[rgb]{0.56,0.35,0.01}{\textbf{\textit{#1}}}}
\newcommand{\ConstantTok}[1]{\textcolor[rgb]{0.56,0.35,0.01}{#1}}
\newcommand{\ControlFlowTok}[1]{\textcolor[rgb]{0.13,0.29,0.53}{\textbf{#1}}}
\newcommand{\DataTypeTok}[1]{\textcolor[rgb]{0.13,0.29,0.53}{#1}}
\newcommand{\DecValTok}[1]{\textcolor[rgb]{0.00,0.00,0.81}{#1}}
\newcommand{\DocumentationTok}[1]{\textcolor[rgb]{0.56,0.35,0.01}{\textbf{\textit{#1}}}}
\newcommand{\ErrorTok}[1]{\textcolor[rgb]{0.64,0.00,0.00}{\textbf{#1}}}
\newcommand{\ExtensionTok}[1]{#1}
\newcommand{\FloatTok}[1]{\textcolor[rgb]{0.00,0.00,0.81}{#1}}
\newcommand{\FunctionTok}[1]{\textcolor[rgb]{0.13,0.29,0.53}{\textbf{#1}}}
\newcommand{\ImportTok}[1]{#1}
\newcommand{\InformationTok}[1]{\textcolor[rgb]{0.56,0.35,0.01}{\textbf{\textit{#1}}}}
\newcommand{\KeywordTok}[1]{\textcolor[rgb]{0.13,0.29,0.53}{\textbf{#1}}}
\newcommand{\NormalTok}[1]{#1}
\newcommand{\OperatorTok}[1]{\textcolor[rgb]{0.81,0.36,0.00}{\textbf{#1}}}
\newcommand{\OtherTok}[1]{\textcolor[rgb]{0.56,0.35,0.01}{#1}}
\newcommand{\PreprocessorTok}[1]{\textcolor[rgb]{0.56,0.35,0.01}{\textit{#1}}}
\newcommand{\RegionMarkerTok}[1]{#1}
\newcommand{\SpecialCharTok}[1]{\textcolor[rgb]{0.81,0.36,0.00}{\textbf{#1}}}
\newcommand{\SpecialStringTok}[1]{\textcolor[rgb]{0.31,0.60,0.02}{#1}}
\newcommand{\StringTok}[1]{\textcolor[rgb]{0.31,0.60,0.02}{#1}}
\newcommand{\VariableTok}[1]{\textcolor[rgb]{0.00,0.00,0.00}{#1}}
\newcommand{\VerbatimStringTok}[1]{\textcolor[rgb]{0.31,0.60,0.02}{#1}}
\newcommand{\WarningTok}[1]{\textcolor[rgb]{0.56,0.35,0.01}{\textbf{\textit{#1}}}}
\usepackage{graphicx}
\makeatletter
\newsavebox\pandoc@box
\newcommand*\pandocbounded[1]{% scales image to fit in text height/width
  \sbox\pandoc@box{#1}%
  \Gscale@div\@tempa{\textheight}{\dimexpr\ht\pandoc@box+\dp\pandoc@box\relax}%
  \Gscale@div\@tempb{\linewidth}{\wd\pandoc@box}%
  \ifdim\@tempb\p@<\@tempa\p@\let\@tempa\@tempb\fi% select the smaller of both
  \ifdim\@tempa\p@<\p@\scalebox{\@tempa}{\usebox\pandoc@box}%
  \else\usebox{\pandoc@box}%
  \fi%
}
% Set default figure placement to htbp
\def\fps@figure{htbp}
\makeatother
\setlength{\emergencystretch}{3em} % prevent overfull lines
\providecommand{\tightlist}{%
  \setlength{\itemsep}{0pt}\setlength{\parskip}{0pt}}
\usepackage{bookmark}
\IfFileExists{xurl.sty}{\usepackage{xurl}}{} % add URL line breaks if available
\urlstyle{same}
\hypersetup{
  pdftitle={Constrained Mixture Designs},
  hidelinks,
  pdfcreator={LaTeX via pandoc}}

\title{Constrained Mixture Designs}
\author{}
\date{\vspace{-2.5em}2026-04-17}

\begin{document}
\maketitle

\subsection{Introduction}\label{introduction}

In many real-world mixture experiments, each component carries
individual lower and upper bounds --- salt must remain at least 5 \%,
alcohol must not exceed 15 \%, an active pharmaceutical ingredient must
be between 2 and 8 \% --- so the feasible blend region is no longer the
full simplex but a constrained convex polytope inside it. Standard
simplex-lattice and simplex-centroid designs do not exploit this
structure and would place runs outside the feasible region.
Extreme-vertices designs (McLean and Anderson, 1966) instead enumerate
the vertices of the constraint polytope and sample them, augmenting with
edge midpoints, face centroids, and an overall centroid as the design
budget allows. The result is a design that is simultaneously feasible
everywhere and informative about the response across the full
constrained region.

\subsection{Prerequisites}\label{prerequisites}

A working understanding of simplex-lattice and simplex-centroid mixture
designs, the Scheffé canonical polynomial, and basic convex-geometry
concepts (vertices, edges, polytopes).

\subsection{Theory}\label{theory}

Given lower bounds \(L_i\) and upper bounds \(U_i\) on each of the \(q\)
components, the feasible region is the convex polytope
\(\{x : L_i \leq x_i \leq U_i, \sum_i x_i = 1\}\). Its vertices are
computable from the constraint set and form the natural design points;
the design is augmented with edge midpoints, face centroids, and the
overall centroid as needed for fitting the chosen Scheffé polynomial
order.

The pseudo-component transformation
\(x'_i = (x_i - L_i)/(1 - \sum_i L_i)\) maps the constrained region back
to a standard unit simplex when only lower bounds are active, allowing
standard simplex-lattice or simplex-centroid construction in the
transformed coordinates. Predictions can then be back-transformed to the
original-component scale.

\subsection{Assumptions}\label{assumptions}

Components are non-negative and sum to one (they are proportions of a
fixed total amount), the lower and upper bounds are fixed and known in
advance, and the residual error is approximately Normal with constant
variance over the feasible region.

\subsection{R Implementation}\label{r-implementation}

\begin{Shaded}
\begin{Highlighting}[]
\FunctionTok{library}\NormalTok{(mixexp)}

\NormalTok{lb }\OtherTok{\textless{}{-}} \FunctionTok{c}\NormalTok{(}\FloatTok{0.2}\NormalTok{, }\FloatTok{0.1}\NormalTok{, }\FloatTok{0.3}\NormalTok{)}
\NormalTok{ub }\OtherTok{\textless{}{-}} \FunctionTok{c}\NormalTok{(}\FloatTok{0.5}\NormalTok{, }\FloatTok{0.6}\NormalTok{, }\FloatTok{0.7}\NormalTok{)}

\NormalTok{xv }\OtherTok{\textless{}{-}} \FunctionTok{Xvert}\NormalTok{(}\AttributeTok{nfac =} \DecValTok{3}\NormalTok{, }\AttributeTok{lc =}\NormalTok{ lb, }\AttributeTok{uc =}\NormalTok{ ub, }\AttributeTok{ndm =} \DecValTok{2}\NormalTok{)}
\NormalTok{xv}

\FunctionTok{set.seed}\NormalTok{(}\DecValTok{2026}\NormalTok{)}
\NormalTok{xv}\SpecialCharTok{$}\NormalTok{y }\OtherTok{\textless{}{-}} \DecValTok{10} \SpecialCharTok{*}\NormalTok{ xv}\SpecialCharTok{$}\NormalTok{x1 }\SpecialCharTok{+} \DecValTok{12} \SpecialCharTok{*}\NormalTok{ xv}\SpecialCharTok{$}\NormalTok{x2 }\SpecialCharTok{+} \DecValTok{8} \SpecialCharTok{*}\NormalTok{ xv}\SpecialCharTok{$}\NormalTok{x3 }\SpecialCharTok{+}
        \DecValTok{6} \SpecialCharTok{*}\NormalTok{ xv}\SpecialCharTok{$}\NormalTok{x1 }\SpecialCharTok{*}\NormalTok{ xv}\SpecialCharTok{$}\NormalTok{x2 }\SpecialCharTok{+}
        \FunctionTok{rnorm}\NormalTok{(}\FunctionTok{nrow}\NormalTok{(xv), }\DecValTok{0}\NormalTok{, }\FloatTok{0.3}\NormalTok{)}

\NormalTok{fit }\OtherTok{\textless{}{-}} \FunctionTok{MixModel}\NormalTok{(xv, }\StringTok{"y"}\NormalTok{,}
                \AttributeTok{mixcomps =} \FunctionTok{c}\NormalTok{(}\StringTok{"x1"}\NormalTok{, }\StringTok{"x2"}\NormalTok{, }\StringTok{"x3"}\NormalTok{), }\AttributeTok{model =} \DecValTok{2}\NormalTok{)}
\FunctionTok{summary}\NormalTok{(fit)}
\end{Highlighting}
\end{Shaded}

\subsection{Output \& Results}\label{output-results}

\texttt{Xvert()} returns the vertices of the constrained polytope along
with optionally requested edge midpoints and centroids.
\texttt{MixModel()} fits the Scheffé linear, quadratic, or special-cubic
mixture model on the design, with coefficients interpretable as blend
effects. Visualising the polytope and overlaying the design points
before running the experiment is a recommended sanity check.

\subsection{Interpretation}\label{interpretation}

A reporting sentence: ``The extreme-vertices design identified seven
feasible vertices of the constrained simplex with bounds
\(L = (0.2, 0.1, 0.3)\) and \(U = (0.5, 0.6, 0.7)\). The fitted
quadratic Scheffé model showed the response was maximised near the
vertex \((0.5, 0.2, 0.3)\) with a meaningful pairwise synergy between
components 1 and 2 (\(\beta_{12} = 5.9\), \(p = 0.01\)). Linear blend
coefficients indicated component 2 was individually most effective.''
Always describe the feasible region and report which vertex achieved the
optimum.

\subsection{Practical Tips}\label{practical-tips}

\begin{itemize}
\tightlist
\item
  Use extreme-vertices designs whenever component bounds are genuinely
  binding; interior-only designs (simplex-lattice or centroid) miss
  optima that lie on the boundary of the feasible region --- a common
  occurrence when constraints reflect formulation requirements.
\item
  D-optimal mixture designs (via \texttt{AlgDesign::optMonteCarlo} or
  \texttt{AlgDesign::optFederov}) handle complex constraints elegantly
  and give the most efficient design when the number of feasible
  candidate points exceeds the run budget.
\item
  Pseudo-component analysis makes the fit and interpretation more
  intuitive when only lower bounds are active: fit on the transformed
  unit simplex, then back-transform predictions to original component
  proportions for reporting.
\item
  For ratio constraints (e.g., solvent₁ : solvent₂ = 2 : 1), use a ratio
  reparameterisation rather than naive component-level bounds; the
  constraint algebra is cleaner and the design more efficient.
\item
  Always plot the feasible region in 3D or on the simplex before
  finalising the design; visualising the polytope helps detect
  over-tight constraints (a near-empty polytope) or accidentally
  redundant ones, and reveals which vertices are clinically meaningful.
\item
  When the polytope is degenerate (very thin in one direction), the
  design will be ill-conditioned; relax the offending constraint or
  accept a higher prediction variance in the squeezed direction.
\end{itemize}

\subsection{R Packages Used}\label{r-packages-used}

\texttt{mixexp::Xvert()} for the canonical extreme-vertices design and
\texttt{MixModel()} for Scheffé-model fitting;
\texttt{AlgDesign::optFederov()} and \texttt{optMonteCarlo()} for
D-optimal mixture designs under arbitrary linear constraints;
\texttt{qualityTools::mixDesign()} for an alternative interface;
\texttt{daewr} for textbook companion datasets and worked
constrained-mixture examples; \texttt{MixExpR} for a tidyverse-friendly
mixture-design workflow.

\subsection{For Reviewers}\label{for-reviewers}

\textbf{What to look for in a paper using this method.}

\begin{itemize}
\tightlist
\item
  \textbf{Common misapplications.}
\item
  \textbf{Diagnostics that should be reported but often aren't.}
\item
  \textbf{Red flags in tables and figures.}
\item
  \textbf{What to verify.}
\item
  \textbf{What an adequate Methods paragraph must contain.}
\end{itemize}

\subsection{See also --- labs in this
chapter}\label{see-also-labs-in-this-chapter}

\begin{itemize}
\tightlist
\item
  \href{labs/lab-observational-designs-and-strobe.Rmd}{Observational
  designs and STROBE}
\item
  \href{labs/lab-randomised-controlled-trials.Rmd}{Randomised controlled
  trials}
\item
  \href{labs/lab-adaptive-non-inferiority-equivalence-trials.Rmd}{Adaptive,
  non-inferiority, equivalence trials}
\item
  \href{labs/lab-bench-and-translational-design.Rmd}{Bench and
  translational design}
\item
  \href{labs/lab-mcar-mar-mnar.Rmd}{MCAR, MAR, MNAR}
\item
  \href{labs/lab-multiple-imputation-with-mice.Rmd}{Multiple imputation
  with mice}
\item
  \href{labs/lab-linear-mixed-models-with-lme4.Rmd}{Linear mixed models
  with lme4}
\item
  \href{labs/lab-glmms-and-gee.Rmd}{GLMMs and GEE}
\item
  \href{labs/lab-dags-with-dagitty-and-ggdag.Rmd}{DAGs with dagitty and
  ggdag}
\item
  \href{labs/lab-propensity-scores-and-iptw.Rmd}{Propensity scores and
  IPTW}
\item
  \href{labs/lab-g-methods-iv-did-rdd-heterogeneity-of-treatment-effect.Rmd}{G-methods,
  IV, DiD, RDD; heterogeneity of treatment effect}
\item
  \href{labs/lab-sir-and-seir-models-with-desolve.Rmd}{SIR and SEIR
  models with deSolve}
\item
  \href{labs/lab-pre-registration-and-statistical-analysis-plans.Rmd}{Pre-registration
  and statistical analysis plans}
\end{itemize}

\end{document}
